\section{Mở đầu}
\label{sec:intro}

\subsection{Bối cảnh và động cơ}
Kubernetes đã trở thành nền tảng điều phối container phổ biến nhất cho điện toán
đám mây hiện đại. Khác với máy ảo, các container trên cùng một nút (node)
\emph{chia sẻ chung một nhân} hệ điều hành; ranh giới cô lập giữa chúng vì thế mỏng
hơn nhiều và hội tụ tại một điểm duy nhất: giao diện lời gọi hệ thống
(syscall)---ranh giới giữa không gian người dùng và nhân~\cite{ref1_jarkas_survey,
ref3_sultan_survey}. Một nhân Linux phơi bày hàng trăm syscall, phần lớn không cần
thiết cho một khối lượng công việc cụ thể, khiến bề mặt tấn công phình
to~\cite{ref11_shu_security}.

Hệ quả là \emph{leo thang đặc quyền} (privilege escalation) và \emph{thoát
container} (container escape) trở thành lớp mối đe dọa nguy hiểm nhất: một container
bị xâm phạm có thể lạm dụng syscall đặc quyền để chiếm quyền \texttt{root} trên nút
và phá vỡ cô lập~\cite{ref5_combe_docker}. Các họ lỗ hổng tiêu biểu---PwnKit
(\mbox{CVE-2021-4034})~\cite{cve_2021_4034}, Dirty~Pipe
(\mbox{CVE-2022-0847})~\cite{cve_2022_0847}, thoát \texttt{runc}
(\mbox{CVE-2019-5736})~\cite{cve_2019_5736}---đều để lại \emph{dấu vết hành vi} ở
tầng syscall (gọi \texttt{ptrace}, \texttt{setuid}/\texttt{capset},
\texttt{unshare}/\texttt{setns}, \texttt{mount}, nạp mô-đun nhân, \texttt{bpf}\dots).
Các cơ chế tĩnh như hồ sơ seccomp khó duy trì, dễ chặn nhầm thao tác hợp lệ hoặc bỏ
sót tấn công nhiều bước, và không nhận thức được \emph{chuỗi} hành vi đặc trưng cho
leo thang đặc quyền.

Vì vậy, một hướng tiếp cận thích nghi---\emph{phát hiện bất thường} trên chuỗi
syscall thời gian thực---ngày càng được quan tâm. Framework
DeSFAM~\cite{desfam_original} là một đại diện tiêu biểu: nó kết hợp lập hồ sơ syscall
lai, phát hiện bất thường bằng học máy độ trễ thấp, và thực thi trong nhân qua
eBPF/LSM. Trái tim phát hiện của DeSFAM là mô-đun \emph{SyscallAD}---đối tượng
nghiên cứu của báo cáo này.

\subsection{Đối tượng và phạm vi nghiên cứu}
\label{subsec:scope}
\textbf{Đối tượng nghiên cứu} của báo cáo là \emph{mô-đun phát hiện bất thường
SyscallAD} (SysAD) trong framework DeSFAM. Framework DeSFAM gồm \emph{ba} mô-đun
(giai đoạn) tích hợp:
\begin{enumerate}[label=\textbf{Mô-đun \arabic*.},leftmargin=2.6em]
  \item \textbf{Lập danh sách truy cập syscall lai và sinh hồ sơ seccomp} (Giai
        đoạn~1): tĩnh + động + lọc theo CVE.
  \item \textbf{SyscallAD---giám sát thời gian thực và phát hiện bất thường} (Giai
        đoạn~2): tập hợp VAE\,+\,iForest trên chuỗi syscall theo cửa sổ trượt.
  \item \textbf{Thực thi thích ứng và quản lý chính sách} (Giai đoạn~3): chấm điểm
        rủi ro, ánh xạ MITRE ATT\&CK, chặn trong nhân.
\end{enumerate}

\textbf{Phạm vi}: báo cáo \emph{tập trung} vào Mô-đun~2 (SyscallAD)---cơ chế phát
hiện leo thang đặc quyền dựa trên hành vi syscall. Mô-đun~1 và Mô-đun~3 \emph{chỉ
được điểm qua} (Mục~\ref{sec:desfam}) nhằm định vị SyscallAD trong tổng thể kiến
trúc, \emph{không} phải đối tượng nghiên cứu chính và không được hiện thực hóa lại.
Ngữ cảnh ứng dụng được chọn là phát hiện leo thang đặc quyền trong Kubernetes.

\subsection{Mục tiêu nghiên cứu}
Là một báo cáo học thuật xây dựng trên một công trình đã công bố (DeSFAM), mục tiêu
của báo cáo \emph{không} phải đề xuất một phương pháp mới mà là \emph{nghiên cứu và
tái lập} mô-đun SyscallAD nhằm:
\begin{enumerate}[label=(\roman*),leftmargin=2.0em]
  \item hiểu và hệ thống hóa thiết kế của SyscallAD---cơ chế biến chuỗi syscall
        thành đặc trưng, mô hình học máy, hợp nhất điểm và chọn ngưỡng---cùng cách nó
        nắm bắt dấu vết hành vi của leo thang đặc quyền;
  \item \emph{tái lập trung thực} SyscallAD và chạy lại thực nghiệm trên DongTing để
        \textbf{cung cấp bằng chứng sơ bộ về tính khả thi} (feasibility)---trong phạm
        vi bộ dữ liệu và kịch bản \emph{mô phỏng} đã dùng---rằng cơ chế phát hiện vận
        hành \emph{phù hợp} với công bố, với quy trình đóng gói bằng Docker có thể
        chạy lại độc lập. Báo cáo phân biệt rõ ``\emph{tái lập trung thành với nguồn}''
        (cùng pipeline, cùng bậc số liệu) với ``\emph{chứng minh cơ chế đúng}'': cái
        sau cần thí nghiệm cắt bỏ ở Mục~\ref{subsec:confound} và do đó chỉ được kết
        luận thận trọng;
  \item tạo một nền tảng tái lập được, làm cơ sở cho các nghiên cứu mở rộng về sau
        (bổ sung đặc trưng, mở rộng khối lượng công việc, kết nối khâu thực thi).
\end{enumerate}

\subsection{Đóng góp}
So với việc chỉ \emph{lặp lại} DeSFAM, báo cáo có ba đóng góp cụ thể:
\begin{enumerate}[label=(\arabic*),leftmargin=2.0em]
  \item một bản \emph{tái hiện độc lập} mô-đun SyscallAD kèm bộ tạo phẩm đóng gói
        bằng Docker, chạy lại được;
  \item một \emph{bộ công cụ đánh giá nghiêm ngặt} (đóng gói trong
        \texttt{Experiment/}): chỉ số mức cửa sổ \emph{và} mức chuỗi, thí nghiệm cắt
        bỏ để kiểm tra nhiễu thiết kế, kiểm toán mã hóa, và tiêu chí đánh giá đặt
        trước---những thứ đường ống gốc thiếu.
\end{enumerate}

\noindent\textbf{Phạm vi báo cáo này:} hoàn tất \emph{đường ống ngoại tuyến}
(huấn luyện $\to$ kiểm tra $\to$ đánh giá) của mô-đun SyscallAD trên DongTing. Việc
phục vụ/thực thi trực tiếp (triển khai trên hạ tầng container) \emph{không} thuộc
phạm vi báo cáo.

\paragraph{Giới hạn mô hình đe dọa.} Báo cáo chỉ xét leo thang đặc quyền/thoát
container \emph{quan sát được ở tầng syscall, mức một nút}. Các con đường leo thang
ở \emph{tầng điều phối} (đánh cắp token ServiceAccount, RBAC quá rộng, lạm dụng dịch
vụ metadata) phần lớn \emph{vô hình} ở tầng syscall mức nút và nằm ngoài phạm vi.

\subsection{Bố cục}
Báo cáo (i) hệ thống hóa thiết kế mô-đun SyscallAD và làm rõ cơ chế phát hiện leo
thang đặc quyền của nó; (ii) tái lập SyscallAD và đánh giá thực nghiệm trên DongTing;
(iii) đối chiếu \emph{trung thực} với số liệu công bố và rút ra bài học triển khai.

Phần còn lại được tổ chức như sau: Mục~\ref{sec:background} trình bày cơ sở lý
thuyết; Mục~\ref{sec:desfam} tổng quan framework DeSFAM và điểm qua Mô-đun~1, 3;
Mục~\ref{sec:syscallad} đi sâu vào kiến trúc và thuật toán của SyscallAD;
Mục~\ref{sec:privesc} phân tích cách SyscallAD phát hiện leo thang đặc quyền;
Mục~\ref{sec:setup} mô tả thiết lập thực nghiệm; Mục~\ref{sec:results} báo cáo kết
quả quan sát; Mục~\ref{sec:discussion} thảo luận; Mục~\ref{sec:conclusion} kết luận.
